%% Tikz settings for Feynman diagrams

\definecolor{betterYellow}{HTML}{FFBF00}
\usetikzlibrary{decorations.pathreplacing,decorations.markings,decorations.pathmorphing,arrows.meta,math,shapes}
\tikzset{
  % style to apply some styles to each segment of a path
  on each segment/.style={
    decorate,
    decoration={
      show path construction,
      moveto code={},
      lineto code={
        \path [#1]
        (\tikzinputsegmentfirst) -- (\tikzinputsegmentlast);
      },
      curveto code={
        \path [#1] (\tikzinputsegmentfirst)
        .. controls
        (\tikzinputsegmentsupporta) and (\tikzinputsegmentsupportb)
        ..
        (\tikzinputsegmentlast);
      },
      closepath code={
        \path [#1]
        (\tikzinputsegmentfirst) -- (\tikzinputsegmentlast);
      },
    },
  },
  % style to add an arrow in the middle of a path as fermion line
  fermion/.style={
    thick,
    draw=#1, % assicura sempre la linea visibile
    postaction={
      on each segment={
        decorate,
        decoration={
          markings,
          mark=at position 0.5 with {\arrow{stealth}}
        }
      }
    }
  },
  photon/.style={
    on each segment = {
        thick,
        draw=#1,
        decorate,
        decoration={wave,amplitude=6pt}
    }
  },
  gluon/.style={
    on each segment = {
        thick,
        draw=#1,
        decorate,
        decoration={gluoncoil,aspect=1,segment length=6pt}
    }
  },
  wzboson/.style={
    on each segment = {
        thick,
        draw=#1,
        decorate,
        decoration={wave,amplitude=6pt}
    }
  },
}

\makeatletter
\pgfdeclaredecoration{gluoncoil}{initial}
{
  \state{initial}[
    width=+0pt,
    next state=coil,
    persistent precomputation={
      \pgfmathsetmacro\matchinglength{
        (ceil(\pgfdecoratedinputsegmentlength / \pgfdecorationsegmentlength) - \pgfdecoratedinputsegmentlength / \pgfdecorationsegmentlength) > 0.5
        ? (\pgfdecoratedinputsegmentlength - 2 * \pgfdecorationsegmentaspect * \pgfdecorationsegmentamplitude) / (floor(\pgfdecoratedinputsegmentlength / \pgfdecorationsegmentlength) + 0.499)
        : (\pgfdecoratedinputsegmentlength - 2 * \pgfdecorationsegmentaspect * \pgfdecorationsegmentamplitude) / (ceil(\pgfdecoratedinputsegmentlength / \pgfdecorationsegmentlength) + 0.499)
      }
      \setlength{\pgfdecorationsegmentlength}{\matchinglength pt}
    },
  ]{}
  \state{coil}[
    switch if less than=%
      0.5\pgfdecorationsegmentlength%
      +\pgfdecorationsegmentaspect\pgfdecorationsegmentamplitude%
      +\pgfdecorationsegmentaspect\pgfdecorationsegmentamplitude to last,
    width=+\pgfdecorationsegmentlength,
  ]
  {
    \pgfpathcurveto
    {\pgfpoint@oncoil{0    }{ 0.555}{ 1}}
    {\pgfpoint@oncoil{0.445}{ 1    }{ 2}}
    {\pgfpoint@oncoil{1    }{ 1    }{ 3}}
    \pgfpathcurveto
    {\pgfpoint@oncoil{1.555}{ 1    }{ 4}}
    {\pgfpoint@oncoil{2    }{ 0.555}{ 5}}
    {\pgfpoint@oncoil{2    }{ 0    }{ 6}}
    \pgfpathcurveto
    {\pgfpoint@oncoil{2    }{-0.555}{ 7}}
    {\pgfpoint@oncoil{1.555}{-1    }{ 8}}
    {\pgfpoint@oncoil{1    }{-1    }{ 9}}
    \pgfpathcurveto
    {\pgfpoint@oncoil{0.445}{-1    }{10}}
    {\pgfpoint@oncoil{0    }{-0.555}{11}}
    {\pgfpoint@oncoil{0    }{ 0    }{12}}
  }
  \state{last}[next state=final]
  {
    \pgfpathcurveto
    {\pgfpoint@oncoil{0    }{ 0.555}{1}}
    {\pgfpoint@oncoil{0.445}{ 1    }{2}}
    {\pgfpoint@oncoil{1    }{ 1    }{3}}
    \pgfpathcurveto
    {\pgfpoint@oncoil{1.555}{ 1    }{4}}
    {\pgfpoint@oncoil{2    }{ 0.555}{5}}
    {\pgfpoint@oncoil{2    }{ 0    }{6}}
  }
  \state{final}{}
}
\newif\ifstartcompletesineup
\newif\ifendcompletesineup
\pgfkeys{
    /pgf/decoration/.cd,
    start up/.is if=startcompletesineup,
    start up=true,
    start up/.default=true,
    start down/.style={/pgf/decoration/start up=false},
    end up/.is if=endcompletesineup,
    end up=true,
    end up/.default=true,
    end down/.style={/pgf/decoration/end up=false}
}
\pgfdeclaredecoration{wave}{initial}
{
    \state{initial}[
        width=+0pt,
        next state=upsine,
        persistent precomputation={
            \ifstartcompletesineup
                \pgfkeys{/pgf/decoration automaton/next state=upsine}
                \ifendcompletesineup
                    \pgfmathsetmacro\matchinglength{
                        0.5*\pgfdecoratedinputsegmentlength / (ceil(0.5* \pgfdecoratedinputsegmentlength / \pgfdecorationsegmentlength) )
                    }
                \else
                    \pgfmathsetmacro\matchinglength{
                        0.5 * \pgfdecoratedinputsegmentlength / (ceil(0.5 * \pgfdecoratedinputsegmentlength / \pgfdecorationsegmentlength ) - 0.499)
                    }
                \fi
            \else
                \pgfkeys{/pgf/decoration automaton/next state=downsine}
                \ifendcompletesineup
                    \pgfmathsetmacro\matchinglength{
                        0.5* \pgfdecoratedinputsegmentlength / (ceil(0.5 * \pgfdecoratedinputsegmentlength / \pgfdecorationsegmentlength ) - 0.4999)
                    }
                \else
                    \pgfmathsetmacro\matchinglength{
                        0.5 * \pgfdecoratedinputsegmentlength / (ceil(0.5 * \pgfdecoratedinputsegmentlength / \pgfdecorationsegmentlength ) )
                    }
                \fi
            \fi
            \setlength{\pgfdecorationsegmentlength}{\matchinglength pt}
        }] {}
    \state{downsine}[width=\pgfdecorationsegmentlength,next state=upsine]{
        \pgfpathsine{\pgfpoint{0.5\pgfdecorationsegmentlength}{0.5\pgfdecorationsegmentamplitude}}
        \pgfpathcosine{\pgfpoint{0.5\pgfdecorationsegmentlength}{-0.5\pgfdecorationsegmentamplitude}}
    }
    \state{upsine}[width=\pgfdecorationsegmentlength,next state=downsine]{
        \pgfpathsine{\pgfpoint{0.5\pgfdecorationsegmentlength}{-0.5\pgfdecorationsegmentamplitude}}
        \pgfpathcosine{\pgfpoint{0.5\pgfdecorationsegmentlength}{0.5\pgfdecorationsegmentamplitude}}
}
    \state{final}{}
}

\makeatother
\newcommand{\addfeynman}[3][]{
    \begin{figure}[H]  % !h or H
        \centering
        \tikz[scale=#3,transform shape,every node/.style={font=\footnotesize}]{
            #2
        }
        \caption{#1}
        \label{feynm:#1}
    \end{figure}
}
\newcommand{\addfeynmantikz}[2]{
        \tikz[scale=#2,transform shape,every node/.style={font=\footnotesize}]{
            #1
        }
}
\newcommand{\addtikz}[4][]{
    \begin{figure}[H]  % !h or H
        \centering
        \tikz[scale=#3,transform shape,every node/.style={font=\footnotesize}]{
            #2
        }
        \caption{#1}
        \label{tikz:#4}
    \end{figure}
}
